\documentclass[aps,prl,twocolumn]{revtex4-2}
\usepackage{amsmath}
\usepackage{physics}

\begin{document}

\title{A short note on biased-noise thresholds}
\author{A. Researcher}
\affiliation{Department of Physics, Example University}

\begin{abstract}
We show that tailoring the checks of a planar code to dephasing noise raises its threshold, and we give a simple argument for why.
\end{abstract}

\maketitle

Biased noise is common in superconducting qubits~\cite{tuckett2018}. Here we give a short argument that a planar code adapted to such noise has a higher threshold than the standard construction.

Consider a planar lattice whose plaquette checks are rotated so that pure dephasing errors map to a repetition code. We find that the effective distance against such errors doubles, which explains the higher threshold reported in numerical studies~\cite{tuckett2018}.

This simple picture suggests that modest changes to a code can pay off on real hardware.

\begin{acknowledgments}
We thank colleagues for discussions.
\end{acknowledgments}

\bibliography{references}

\end{document}
